%=============================================================
% Chapitre 1 — Cadre géographique et contexte électoral
%=============================================================
\chapter{Cadre géographique et contexte électoral}
\label{ch:geo}

\section{La Nouvelle-Aquitaine en chiffres}

La \textbf{Nouvelle-Aquitaine} est la plus grande région métropolitaine française par
superficie, avec environ \textbf{84 000 km²}. Elle regroupe 12 départements issus de la
fusion en 2016 des anciennes régions Aquitaine, Limousin et Poitou-Charentes.

\begin{table}[H]
\centering
\caption{Les 12 départements de Nouvelle-Aquitaine — caractéristiques clés}
\label{tab:departements}
\begin{tabular}{llrr}
\toprule
\textbf{Code} & \textbf{Département} & \textbf{Pop. (2022, k)} & \textbf{Cantons} \\
\midrule
16 & Charente           & 353  & 29 \\
17 & Charente-Maritime  & 650  & 27 \\
19 & Corrèze            & 242  & 19 \\
23 & Creuse             & 115  & 16 \\
24 & Dordogne           & 416  & 25 \\
33 & Gironde            & 1 622 & 43 \\
40 & Landes             & 415  & 16 \\
47 & Lot-et-Garonne     & 335  & 17 \\
64 & Pyrénées-Atlantiques & 693 & 34 \\
79 & Deux-Sèvres        & 372  & 18 \\
86 & Vienne             & 440  & 28 \\
87 & Haute-Vienne       & 372  & 21 \\
\midrule
\multicolumn{2}{l}{\textbf{Total}} & \textbf{6 025} & \textbf{293} \\
\bottomrule
\end{tabular}
\end{table}

\section{Contexte socio-économique régional}

La région présente des dynamiques contrastées :

\begin{itemize}
  \item \textbf{Gironde} (chef-lieu Bordeaux) : moteur économique régional, forte
    croissance démographique, marché immobilier sous tension.
  \item \textbf{Limousin} (Creuse, Corrèze, Haute-Vienne) : vieillissement démographique,
    déclin industriel, rural profond.
  \item \textbf{Charente, Lot-et-Garonne} : économies agricoles avec des poches
    de fragilité économique, électoralement plus incertains.
  \item \textbf{Pays Basque} (Pyrénées-Atlantiques) : dynamisme touristique et
    identité régionale forte, profil électoral atypique.
\end{itemize}

Ces disparités territoriales font de la Nouvelle-Aquitaine un terrain idéal pour
tester la capacité des modèles ML à capturer des signaux hétérogènes.

\section{Résultats réels --- Présidentielle 2022 (1er tour)}

Lors du premier tour de l'élection présidentielle du 10 avril 2022, les résultats
régionaux placent \textbf{Emmanuel Macron} en tête dans \textbf{11 des 12 départements}
de Nouvelle-Aquitaine. Le Lot-et-Garonne constitue l'exception notable, avec une
victoire de \textbf{Marine Le Pen} (RN).

\begin{table}[H]
\centering
\caption{Résultats réels du 1\textsuperscript{er} tour — présidentielle 2022 en Nouvelle-Aquitaine}
\label{tab:resultats_reels}
\begin{tabular}{lll}
\toprule
\textbf{Département} & \textbf{1\textsuperscript{er} (réel)} & \textbf{2\textsuperscript{e} (réel)} \\
\midrule
Charente            & Macron  & Mélenchon \\
Charente-Maritime   & Macron  & Mélenchon \\
Corrèze             & Macron  & Mélenchon \\
Creuse              & Macron  & Le Pen \\
Deux-Sèvres         & Macron  & Mélenchon \\
Dordogne            & Macron  & Mélenchon \\
Gironde             & Macron  & Mélenchon \\
Haute-Vienne        & Macron  & Mélenchon \\
Landes              & Macron  & Mélenchon \\
\textbf{Lot-et-Garonne} & \textbf{Le Pen} & Macron \\
Pyrénées-Atlantiques & Macron & Mélenchon \\
Vienne              & Macron  & Mélenchon \\
\bottomrule
\end{tabular}
\end{table}

\section{Variable cible et modèle d'orientation électorale}

La variable cible du projet est \textbf{synthétique} et à 5 classes, reflétant la
dynamique socio-économique d'un canton :

\begin{table}[H]
\centering
\caption{Variable cible — 5 classes et orientation électorale associée}
\label{tab:classes}
\begin{tabular}{llll}
\toprule
\textbf{Classe} & \textbf{Signification} & \textbf{Proportion} & \textbf{Orientation} \\
\midrule
Boom        & Croissance forte  & 10\% & Extrême-Gauche \\
Croissance  & Croissance modérée & 20\% & Gauche \\
Stable      & Stabilité         & 40\% & Centre \\
Déclin      & Déclin modéré     & 20\% & Droite \\
Crise       & Crise économique  & 10\% & Extrême-Droite \\
\bottomrule
\end{tabular}
\end{table}

\noindent La correspondance entre classe économique et orientation (\texttt{ECO\_TO\_ORIENTATION})
permet, à partir de la distribution des prédictions sur les cantons d'un département,
de calculer la probabilité de vote pour chacun des 12 candidats du 1\textsuperscript{er} tour.

Le MCD comprend les entités principales : \texttt{Canton}, \texttt{Indicateur2016},
\texttt{Indicateur2022}, \texttt{DeltaFeature}, \texttt{PredictionDept} et
\texttt{Candidat}, reliées par des associations de type 1-N et N-M.
